%%%%%%%%%%%%%%%%%%%%%%%%%%%%%%%%%%%%%%%%%%%%%%%%%%%%%%%%%%%%%%%%%%%%%%%%
% Plantilla TFG/TFM
% Universidad de Murcia. Facultad de Informática
%%%%%%%%%%%%%%%%%%%%%%%%%%%%%%%%%%%%%%%%%%%%%%%%%%%%%%%%%%%%%%%%%%%%%%%%


\chapter*{Resumen}

En la era digital actual, el acceso eficiente y preciso a grandes volúmenes de información estructurada y no estructurada representa uno de los mayores desafíos en el ámbito de las Ciencias de la Computación. Históricamente, la Web Semántica y los Grafos de Conocimiento (\textit{Knowledge Graphs}) han proporcionado un marco tecnológico robusto para modelar información de manera determinista mediante tripletas (Sujeto, Predicado, Objeto) basadas en el estándar RDF (Resource Description Framework). El acceso a esta información altamente estructurada y veraz se realiza tradicionalmente a través de SPARQL, un lenguaje de consulta potente pero con una curva de aprendizaje pronunciada que supone una barrera técnica insalvable para la mayoría de los usuarios no expertos.

De forma paralela, el reciente auge de los Modelos de Lenguaje de Gran Escala (LLMs) ha revolucionado la interacción humano-máquina. Estos modelos permiten establecer interfaces conversacionales fluidas en lenguaje natural, facilitando enormemente el acceso a la información. Sin embargo, su naturaleza estocástica los hace vulnerables a problemas inherentes como las denominadas "alucinaciones" (la generación de respuestas que resultan gramaticalmente correctas y plausibles, pero que son total o parcialmente falsas) y presentan serias dificultades para recuperar datos precisos y actualizados de dominios de conocimiento altamente específicos.

El presente Trabajo de Fin de Grado (TFG) aborda la convergencia de ambas tecnologías con el objetivo de aunar la flexibilidad y capacidad de comprensión del lenguaje natural propia de los LLMs con la precisión factual inmutable de un Grafo de Conocimiento. Para lograr este ambicioso objetivo, se ha diseñado, implementado y evaluado un agente conversacional inteligente (chatbot) especializado en el dominio cinematográfico. A diferencia de un chatbot tradicional que responde basándose puramente en los patrones aprendidos en su entrenamiento, este sistema actúa como una interfaz transparente que traduce la intención del usuario a consultas SPARQL válidas, ejecutándolas contra una base de datos local y devolviendo la respuesta exacta.

La arquitectura del sistema propuesto se fundamenta en una evolución del paradigma de Generación Aumentada por Recuperación (RAG), elevándolo hacia un sistema de Inteligencia Artificial Agéntica. La solución tecnológica se ha estructurado en tres capas principales. En la capa de presentación, se ha desarrollado una interfaz gráfica de usuario amigable mediante Streamlit, que no solo facilita la interacción conversacional, sino que también expone de manera transparente el razonamiento interno del agente ("Chain-of-Thought") y las consultas generadas en tiempo real. En la capa lógica, el núcleo cognitivo del sistema, impulsado por el modelo Llama 3 ejecutado localmente mediante Ollama, se encarga de la extracción de entidades, su desambiguación contra el vocabulario de Wikidata y la generación dinámica de código SPARQL. Finalmente, la capa de datos descansa sobre un Grafo de Conocimiento local alojado en el motor de base de datos Blazegraph.

Un hito técnico fundamental de este trabajo ha sido la construcción de este repositorio de datos semántico. Depender directamente de \textit{endpoints} públicos como el de Wikidata durante la ejecución en tiempo real introducía niveles inaceptables de latencia y riesgo de bloqueos por límites de peticiones (APIs \textit{rate limits}). Por ello, se implementó un proceso de extracción \textit{offline} estructurado para aislar un subgrafo densamente conectado del dominio cinematográfico (abarcando miles de películas, directores, actores, géneros y fechas de estreno), el cual fue serializado e indexado localmente. Esto proporciona al agente un entorno de ejecución local, excepcionalmente rápido y completamente determinista.

El aspecto más innovador de la implementación radica en el mecanismo autónomo de autocorrección del agente, denominado "bucle de reflexión" (\textit{Reflexion Loop}). La traducción de lenguaje natural a sintaxis SPARQL estricta es propensa a errores. Para mitigarlo, el agente cuenta con un diccionario de propiedades permitidas altamente restrictivo (por ejemplo, \texttt{wdt:P57} exclusivamente para referenciar a un "director") que acota drásticamente su espacio de búsqueda. Si el modelo genera una consulta con errores de sintaxis o que no compila, la capa de datos captura la excepción y, en lugar de fallar de manera abrupta, realimenta al LLM con este error inyectado en un \textit{prompt} de reflexión. De este modo, el agente analiza lógicamente su propia equivocación y reescribe la consulta desde cero, logrando un notable incremento en la robustez, resiliencia y autonomía del sistema.

Para validar el sistema propuesto de forma objetiva, rigurosa y cuantitativa, se desarrolló ad-hoc un marco de evaluación automatizado. Se confeccionó un conjunto de pruebas (\textit{dataset}) compuesto por 103 consultas complejas en lenguaje natural que abarcaban diversas relaciones semánticas de dificultad incremental. El entorno de evaluación puso a prueba seis modelos LLM locales diferentes (Llama 3, Llama 3.2:1B, Mistral, Gemma 2:9B, Gemma 2:2B y DeepSeek-R1:8B) para analizar su rendimiento de forma comparativa. El modelo seleccionado para la versión final del proyecto, Llama 3, demostró un rendimiento excepcional: alcanzó una tasa de precisión del 97.1\% en la fase crítica de extracción de entidades, un 97.1\% en la generación de código SPARQL válido y ejecutable, y un porcentaje de éxito final del 94.3\% entregando la respuesta correcta a la pregunta planteada.

Además de las tasas de éxito en precisión, se midió exhaustivamente el tiempo de ejecución por pregunta, un factor sumamente crítico para asegurar una buena experiencia de usuario en cualquier asistente conversacional interactivo. Llama 3 demostró un rendimiento óptimo en este aspecto, promediando un tiempo de ejecución de respuesta de tan solo 9.51 segundos por pregunta, erigiéndose de manera destacada como el modelo más equilibrado y rápido de entre todos los evaluados en este trabajo. Un análisis pormenorizado de los escasos fallos residuales reveló que estos se debían, en su inmensa mayoría, a ambigüedades idiomáticas extremas por parte del usuario o a cuellos de botella aislados en las llamadas de desambiguación de entidades, y no a fallos estructurales o lógicos en la formulación de consultas deductivas por parte del LLM.

En conclusión, los resultados empíricos obtenidos a lo largo de las pruebas confirman sin lugar a dudas que la integración de agentes de IA fuertemente tipados con Grafos de Conocimiento locales constituye una solución tecnológica sumamente viable y efectiva para la creación de sistemas de pregunta-respuesta (QA) totalmente transparentes y estructuralmente libres de alucinaciones. El sistema diseñado ha logrado el hito de ocultar la enorme complejidad inherente al modelo de datos de la Web Semántica de cara al usuario final, manteniendo no obstante intacta la rigurosa precisión y veracidad de sus datos subyacentes. El presente trabajo sienta, por tanto, unas bases sólidas y prometedoras para futuras líneas de investigación y desarrollo, tales como la mejora de los algoritmos basados en similitud de vectores (\textit{embeddings}) para lograr una desambiguación semántica libre de fallos, o la ampliación escalonada de la ontología base para llegar a soportar razonamientos deductivos mucho más complejos de múltiples saltos lógicos (\textit{multi-hop reasoning}) que permitan extraer conclusiones de datos aparentemente inconexos.

\vspace{1cm}
\textbf{Palabras clave:} Web Semántica, Grafos de Conocimiento, SPARQL, Modelos de Lenguaje (LLMs), Llama3, Generación Aumentada por Recuperación (RAG), Inteligencia Artificial Agéntica, Streamlit, Blazegraph, Mitigación de Alucinaciones.
